\section{Limitations and threats to validity}
\label{sec:limits}

\paragraph{The benchmark is observational and covers one effect on one data set}
C7 compares a model with a response estimated from observational data: the
treated stores and weeks are inferred, and the 8-week estimate is close to a
placebo at false dates. That C7 can tell a response several times too large from
a correct one is shown in the known-answer study, where the truth is known. On
Dominick's the benchmark is too noisy to establish even that, and comparing the
models with it is descriptive. It measures the response
of a store's category units to a sustained change in its regular prices. It says nothing directly about promotions,
displays or temporary cuts, or about M5, Frat and Favorita, where no such variation is used. A model can be right about regular prices and wrong about
promotions; the known-answer study shows how far the other checks go in
covering the rest.

\paragraph{The deviations are inferred}
The public files do not mark the experimental stores or weeks. Some deviations
can be zone moves made for other reasons, and a store moved for a reason that
also moves its sales would bias $\tau_e$. The event study and the placebo
address this for trends that start before the price moves; they cannot rule out
a shock that starts in the same week.

\paragraph{Few control stores}
Zones have few stores, so each scored episode has \zControlsLo{} to
\zControlsHi{} controls. Intervals clustered
by store with so few clusters are too narrow (\S\ref{sec:zones}); the pooled
slope across episodes carries the comparison.

\paragraph{The baseline is a choice}
Every lift is measured against nothing planned, with the audit's regular price
taken as the highest of the last 12 training periods. A different rule gives
different baselines and lifts. The audit applies its rule the same way to every
model and to the truth in the known-answer study; each model's own regular
price can differ from it.

\paragraph{The versions are not a designed sample}
The audited versions are what the companion paper saved. They share data and
differ in many settings at once, so the correlation across versions is
descriptive; it is computed within data set and tested by permutation within
data set, but versions are not independent draws, and the windows of one
version differ in difficulty, so the result depends on whether version-windows
or versions are the unit. Both are reported.

\paragraph{The audit's definition of nothing planned}
The audit builds ``nothing planned'' from the data: levers off, the price at
the series' highest price over the last 12 training periods, and the planned
price where the series has no price in those periods. The models build their
own from their history window: a recorded or held regular price where the
configuration says so, otherwise the highest price over the last 12 periods,
falling back to the series' mean price over the window, or to 1.0 where it
records no price. C1 under the audit's plan therefore measures the distance
between two definitions as much as the model; under the model's own plan it
passes on every window of the final tuned models (\S\ref{sec:tuned}) and was not
scored for the earlier versions (\S\ref{sec:results}). Neither
definition is the retailer's regular price, which these data do not record.

\paragraph{Refits for C6 and C7}
C1 to C5 and C8 read the saved final models. C6 compares them with LightGBM
refitted on the same windows, and C7 uses world models and LightGBM fitted
again with the final settings up to the week before each episode. These refits
are not the runs whose scores the companion paper reports, and the difference
between them and those runs was not measured.

\paragraph{The generated world is simpler than retail}
The known-answer data are multiplicative in the plan, with one shock process
and no cannibalisation between items. They test whether a check can catch a
failure, not how often the failure occurs in practice. The sensitivity worlds change one setting at a
time around one world, and the decision scores (choosing half the planned weeks, ordering to the stated
lift) are simple stand-ins for a planner's choices.
